

Artificial Intelligence represents the broad ambition to create systems capable of autonomous reasoning and learning. Within this field, Deep Learning, a specialized subset, employs multi-layered neural networks to automatically learn complex representations without manual feature engineering. This capability has produced breakthrough results in computer vision and natural language processing, fueled by massive datasets and advances in computational power \cite{Goodfellow2016}. Yet this success comes at a cost: state-of-the-art models now contain billions of parameters, creating critical deployment challenges. Understanding these architectures deeply is essential not only for building them, but for compressing them effectively which is the central challenge of this report.

This chapter establishes the theoretical foundation for these systems. We begin with the artificial neuron, examining its mathematical formulation and the role of non-linearity in capturing complex patterns. We construct the Multi-Layer Perceptron, analyze its architecture and information flow, and establish its expressive power through the universal approximation theorem \cite{hornik1989}. We formalize the optimization framework through loss functions and gradient descent, then explain backpropagation, the algorithm that computes gradients efficiently \cite{Rumelhart1986}. 

While recurrent architectures provided early solutions for sequence modeling, their sequential nature imposed significant limitations. We therefore examine the transformer, which achieves parallel computation through self-attention \cite{vaswani2017attention}. We dissect its encoder-decoder architecture and attention mechanism, building intuition through the Query-Key-Value framework. We analyze multi-head attention's ability to capture diverse relationships and conclude by examining the quadratic complexity bottleneck, a key target for the compression techniques developed in subsequent chapters.

\newpage
\section{The Neuron as a Linear Model}

\subsection{Mathematical Formulation}
The artificial neuron, formally introduced by \cite{McCulloch1943}, serves as the fundamental building block of neural networks. Mathematically, it performs an \textbf{affine transformation} on its inputs followed by a non-linear activation.

\begin{equation}
z = \vect{w}^T \vect{x} + b = \sum_{i=1}^{n} w_i x_i + b
\end{equation}
\noindent Given an input $\vect{x} \in \reals^n$, weights $\vect{w} \in \reals^n$ scale the importance of each feature, while the bias $b \in \reals$ acts as an offset. This bias allows the decision boundary to shift independent of input values \cite{Goodfellow2016}. The resulting logit $z$ is subsequently passed through a non-linear activation function, $a = \phi(z)$.

The term $\vect{w}^T \vect{x}$ measures the alignment between the input and weight vectors. A large magnitude indicates a strong correlation (or anti-correlation) between the input pattern and the feature the neuron is tuned to detect \cite{Haykin1999}.

\subsubsection{Diagram: Single Neuron with Affine Transformation}
Figure~\ref{fig:single_neuron} illustrates the computational flow: inputs are weighted, summed with the bias, and aggregated into the pre-activation $z$.

\begin{figure}[h]
\centering
\begin{tikzpicture}[
    scale=0.6,
    node distance=2cm,
    input/.style={circle, draw=slateblue, fill=mistyblue!30, thick, minimum size=0.8cm},
    neuron/.style={circle, draw=dustyteal, fill=sagegreen!30, thick, minimum size=1.5cm},
    output/.style={circle, draw=slateblue, fill=mistyblue!30, thick, minimum size=0.8cm},
    arrow/.style={-Stealth, thick, color=deepmutedblue}
  ]

  \node[input] (x1) at (0,2) {$x_1$};
  \node[input] (x2) at (0,0.7) {$x_2$};
  \node[input] (x3) at (0,-0.7) {$x_3$};
  \node at (0,-1.7) {$\vdots$};
  \node[input] (xn) at (0,-2.7) {$x_n$};

  \node[input] (bias) at (0,3.5) {$1$};

  \node[neuron] (neuron) at (4,0.5) {$\Sigma$};
  \node[below=0.1cm of neuron, color=deepmutedblue] {\small $z = \vect{w}^T\vect{x} + b$};

  % output z
  \node[output] (z) at (7,0.5) {$z$};

  % connections arrow with their respective weights
  \draw[arrow] (x1) -- node[above, slateblue, font=\small] {$w_1$} (neuron);
  \draw[arrow] (x2) -- node[above, slateblue, font=\small] {$w_2$} (neuron);
  \draw[arrow] (x3) -- node[above, slateblue, font=\small] {$w_3$} (neuron);
  \draw[arrow] (xn) -- node[above, slateblue, font=\small] {$w_n$} (neuron);
  \draw[arrow] (bias) -- node[above, slateblue, font=\small] {$b$} (neuron);

  % output connection
  \draw[arrow] (neuron) -- (z);

\end{tikzpicture}
\caption{Computational graph of a single artificial neuron. Inputs $x_i$ are weighted by $w_i$ and summed with bias $b$ to produce $z$.}
\label{fig:single_neuron}
\end{figure}
\subsection{Geometric Interpretation: The Hyperplane Decision Boundary}
The equation $\vect{w}^T \vect{x} + b = 0$ defines a hyperplane partitioning the space $\reals^n$ into two regions. This geometric view, central to the Perceptron \cite{Rosenblatt1958}, highlights the roles of the parameters:
\begin{itemize}
\item \textbf{Weight Vector $\vect{w}$:} Defines the orientation of the hyperplane. It is orthogonal to the decision surface.
\item \textbf{Bias $b$:} Determines the distance of the hyperplane from the origin, given by $-b/||\vect{w}||_2$.
\end{itemize}

The pre-activation $z$ is proportional to the signed distance from $\vect{x}$ to the hyperplane. Consequently, a single neuron functions as a binary linear classifier, capable of solving linearly separable problems (e.g., AND, OR) but failing on non-linear tasks such as XOR \cite{Minsky1969}.

\subsubsection{Example: Geometric Interpretation in 2D}
Consider an example in $\reals^2$ with weights $\vect{w} = [1, 2]^T$ and bias $b = -3$. The decision boundary is the line where $z=0$:
\begin{equation}
x_1 + 2x_2 - 3 = 0 \quad \Rightarrow \quad x_2 = \frac{3 - x_1}{2}
\end{equation}

Figure~\ref{fig:geometric_2d} visualizes this boundary. The weight vector is normal to the line, pointing toward the positive half-space ($z > 0$).

\begin{figure}[h]
\centering
\begin{tikzpicture}
  \begin{axis}[
      width=12cm,
      height=10cm,
      axis lines=middle,
      xlabel={$x_1$},
      ylabel={$x_2$},
      xmin=-1, xmax=5,
      ymin=-1, ymax=4,
      xtick={0,1,2,3,4,5},
      ytick={0,1,2,3,4},
      grid=major,
      grid style={dashed, gray!30},
      legend pos=north east,
      every axis x label/.style={at={(ticklabel* cs:1.02)}, anchor=west},
      every axis y label/.style={at={(ticklabel* cs:1.02)}, anchor=south},
    ]

    % Decision boundary line: x_2 = (3 - x_1)/2
    \addplot[thick, color=dustyteal, domain=-1:5, samples=2] {(3-x)/2};
    \addlegendentry{Decision boundary: $z=0$}

    % --- REGION z > 0
    \addplot [
      draw=none,
      fill=sagegreen,
      fill opacity=0.2,
      domain=-1:5,
    ] { (3-x)/2 } |- (axis cs:-1,4) -- cycle;
    \node[color=sagegreen!70!black] at (axis cs:1,3.2) {\Large $z > 0$};

    % --- REGION z < 0
    \addplot [
      draw=none,
      fill=mistyblue,
      fill opacity=0.2,
      domain=-1:5,
    ] { (3-x)/2 } |- (axis cs:-1,-1) -- cycle;
    \node[color=mistyblue!70!black] at (axis cs:1,-0.5) {\Large $z < 0$};

    % --- SAMPLE POINTS ---
    % Green points (z > 0)
    \addplot[only marks, mark=*, mark size=3pt, color=sagegreen!70!black]
    coordinates {(1,2.2) (0.5,2.5) (2,2)};
    % Blue points (z < 0)
    \addplot[only marks, mark=*, mark size=3pt, color=mistyblue!70!black]
    coordinates {(2.5,0) (3,-0.2) (4,-0.8)};

    % --- WEIGHT VECTOR ---
    \draw[-Stealth, very thick, color=slateblue]
    (axis cs:2,0.5) -- (axis cs:2.4,1.3)
    node[right, color=slateblue] {$\vect{w} =
      \begin{bmatrix}1 \\ 2
    \end{bmatrix}$};

    % --- DISTANCE INDICATION ---
    \draw[dashed, thick, color=dustyteal]
    (axis cs:0,0) -- (axis cs:0.6,1.2);
    \node[below right, color=dustyteal, font=\small]
    at (axis cs:0.4,0.9)
    {$\frac{-b}{\|\vect{w}\|_2}$};

  \end{axis}
\end{tikzpicture}
\caption{Geometric interpretation in 2D space. The line $x_1 + 2x_2 - 3 = 0$ divides the plane. The weight vector $\vect{w}$ is normal to the decision boundary, pointing toward the region where $z > 0$.}
\label{fig:geometric_2d}
\end{figure}
